\lampiransection{Checklist Data Penelitian Utama}
\label{app:cek-data-penelitian-utama}

Lampiran ini digunakan sebagai sarana cek dan input data penelitian utama. Setelah observasi GrabFood tanggal 13 Juni 2026 dan observasi kompetitor tanggal 19 dan 20 Juni 2026, sebagian data platform sudah tersedia sebagai temuan awal. Data yang belum tersedia tetap dicatat sebagai kebutuhan pengumpulan data berikutnya agar Bab IV dapat ditriangulasikan dengan wawancara, survei, dan dokumentasi usaha.

\begin{table}[!htbp]
    \centering
    \caption{Checklist Kebutuhan Data Penelitian Utama Setelah Observasi GrabFood (Bagian 1)}
    \label{tab:checklist-kebutuhan-data-utama}
    \footnotesize
    \setlength{\tabcolsep}{3pt}
    \renewcommand{\arraystretch}{0.95}
    \begin{tabularx}{\textwidth}{|Z{2.5cm}|Y|Y|Y|}
        \toprule
        \textbf{Kelompok Data} & \textbf{Data yang Sudah Tersedia} & \textbf{Data yang Perlu Dilengkapi} & \textbf{Catatan Analisis} \\ \hline
        \midrule
        Identitas usaha & Profil GrabFood Nasi Gerilya terlihat melalui nama toko, kategori, status toko, rating 4,7, 4.584 penilaian, dan jarak 6,6 km. & Sejarah usaha, struktur pengelolaan, kapasitas produksi, dan kanal penjualan selain GrabFood. & Data platform mendukung TO-GF-01, tetapi profil usaha tetap perlu dikonfirmasi kepada pemilik. \\ \hline
        Kontribusi GrabFood & Pra-survei menyebut GrabFood menjadi kanal penting dan diperkirakan menyumbang lebih dari 50\% pendapatan. & Bukti penjualan, proporsi omzet GrabFood, dan alasan strategis fokus pada GrabFood. & Perlu dokumentasi agar klaim kontribusi tidak hanya berbasis wawancara awal. \\ \hline
        Penjualan & Belum tersedia dalam bentuk angka penjualan bulanan. & Jumlah pesanan, omzet, tren penjualan, periode promo, dan perubahan sebelum-sesudah strategi. & Data ini dibutuhkan untuk mengaitkan SWOT dengan tujuan peningkatan penjualan. \\ \hline
        AOV dan kunjungan & Belum tersedia dalam bentuk dashboard atau rekap angka. & AOV, kunjungan toko, konversi, pembelian ulang, dan pelanggan yang diaktivasi kembali. & Jika data dashboard tidak tersedia, gunakan dokumentasi internal dan keterangan pemilik sebagai pembanding. \\ \hline
        Pelanggan & Ulasan GrabFood menunjukkan tema positif dan negatif/netral, termasuk rasa, porsi, kemasan, salah item, kelengkapan, dan kualitas makanan. & Wawancara pelanggan, survei pendukung, alasan membeli ulang, alasan tidak mengulas, dan pengalaman pelanggan yang tidak tampak pada review. & Perlu mitigasi \textit{survivorship bias} karena ulasan hanya mewakili pelanggan yang memberi review. \\ \hline
        \bottomrule
    \end{tabularx}
    \tablesource{Diolah peneliti berdasarkan observasi GrabFood tanggal 13 Juni 2026 dan observasi kompetitor tanggal 19 dan 20 Juni 2026.}
\end{table}

\begin{table}[!htbp]
    \centering
    \caption{Checklist Kebutuhan Data Penelitian Utama Setelah Observasi GrabFood (Bagian 2)}
    \label{tab:checklist-kebutuhan-data-utama-lanjutan}
    \footnotesize
    \setlength{\tabcolsep}{3pt}
    \renewcommand{\arraystretch}{0.95}
    \begin{tabularx}{\textwidth}{|Z{2.5cm}|Y|Y|Y|}
        \toprule
        \textbf{Kelompok Data} & \textbf{Data yang Sudah Tersedia} & \textbf{Data yang Perlu Dilengkapi} & \textbf{Catatan Analisis} \\ \hline
        \midrule
        Produk & Menu terlaris, menu tidak tersedia, rentang harga, foto produk, dan tema ulasan produk sudah terlihat pada observasi Nasi Gerilya. & Konfirmasi menu unggulan, menu bermasalah, porsi, standar rasa, dan stok dari pemilik/karyawan. & Data awal mendukung TO-GF-04, TO-GF-05, TO-GF-06, TO-GF-07, dan TO-GF-08. \\ \hline
        Harga dan promo & Nasi Gerilya memiliki rentang harga terlihat Rp1.100--Rp275.000; item tersedia Rp1.100--Rp72.000; paket nasi Padang tersedia Rp27.500--Rp72.000; 19 promo/voucher terlihat. & Biaya promo, margin, strategi penetapan harga, dan persepsi pelanggan terhadap harga/ongkir. & Perlu dibaca bersama kompetitor: Paripurna Rp5.000--Rp66.000, Dendeng Batokok Rp5.000--Rp36.000, Istana Krakatau Rp9.000--Rp250.000, Pondok Gurih Rp6.750--Rp246.750, dan Garuda Rp14.300--Rp265.851. \\ \hline
        Operasional & Ulasan menunjukkan indikasi masalah akurasi pesanan, item kurang lengkap, add-on tidak diberikan, menu tidak tersedia, dan konsistensi kualitas. & Alur kerja admin/dapur, pengelolaan stok, jam ramai, SOP pengecekan pesanan, dan penanganan komplain. & Temuan ulasan perlu dikonfirmasi kepada admin/karyawan agar tidak menjadi asumsi sepihak. \\ \hline
        Kompetitor & Lima kompetitor sudah diobservasi: Restoran Paripurna, Rumah Makan Dendeng Batokok, Rumah Makan Istana Krakatau, RM. Pondok Gurih, dan Restoran Garuda. Data mencakup rating, jumlah penilaian, jarak, promo, harga, menu, foto produk, dan ulasan. & Konfirmasi kompetitor yang benar-benar dibandingkan oleh pemilik dan pelanggan, serta observasi ulang bila muncul kompetitor tambahan. & Data awal mendukung TO-KP-01 sampai TO-KP-07 dan dapat dipakai untuk EFAS. \\ \hline
        Bukti visual & 76 bukti visual Nasi Gerilya dan 209 bukti visual kompetitor sudah masuk lampiran bukti. & Bukti observasi lanjutan pada waktu berbeda untuk triangulasi waktu. & Bukti visual sudah cukup sebagai jejak data awal, tetapi tidak menggantikan wawancara dan dokumentasi. \\ \hline
        \bottomrule
    \end{tabularx}
    \tablesource{Diolah peneliti berdasarkan observasi GrabFood tanggal 13 Juni 2026 dan observasi kompetitor tanggal 19 dan 20 Juni 2026.}
\end{table}

\begin{table}[!htbp]
    \centering
    \caption{Log Pengumpulan Data Penelitian Utama}
    \label{tab:log-pengumpulan-data-utama}
    \footnotesize
    \setlength{\tabcolsep}{3pt}
    \renewcommand{\arraystretch}{0.95}
    \begin{tabularx}{\textwidth}{|Z{2.1cm}|Y|Y|Y|}
        \toprule
        \textbf{Tanggal} & \textbf{Kegiatan} & \textbf{Data yang Diperoleh} & \textbf{Kode / Lokasi Arsip} \\ \hline
        \midrule
        13 Juni 2026 & Observasi GrabFood Nasi Gerilya - Glugur Kota & Profil toko, rating, jumlah penilaian, promo, menu, rentang harga, menu tidak tersedia, dan ulasan pelanggan. & TO-GF-01--TO-GF-08; Lampiran~\ref{app:evidence-observasi-grabfood-20260613} \\ \hline
        19 dan 20 Juni 2026 & Observasi GrabFood kompetitor & Profil, rating, jumlah penilaian, jarak, promo, menu, harga, ketersediaan menu, dan ulasan Restoran Paripurna, Rumah Makan Dendeng Batokok, Rumah Makan Istana Krakatau, RM. Pondok Gurih, dan Restoran Garuda. & TO-KP-01--TO-KP-07; Lampiran~\ref{app:evidence-observasi-grabfood-kompetitor-20260619} \\ \hline
        Juni--Juli 2026 & Wawancara pemilik/pengelola & Diisi setelah wawancara utama. & Pedoman pada Lampiran~\ref{app:pedoman-wawancara-utama} \\ \hline
        Juni--Juli 2026 & Wawancara admin/karyawan dan pelanggan & Diisi setelah wawancara operasional dan pelanggan. & Pedoman pada Lampiran~\ref{app:pedoman-wawancara-utama} \\ \hline
        Juni--Juli 2026 & Survei pelanggan pendukung dan dokumentasi kinerja & Diisi setelah survei dan dokumentasi usaha diterima. & Lampiran~\ref{app:survei-pelanggan-pendukung} dan rekap dokumentasi kinerja \\ \hline
        \bottomrule
    \end{tabularx}
    \tablesource{Diolah peneliti sebagai log pengumpulan data penelitian utama (2026).}
\end{table}
